% ============================================================
%  RAPPORT D'ÉCARTS — Phase 4 : Maîtrise (déclenché si écart)
%  Time'Eats — Cellule PMO
% ============================================================
\documentclass[11pt,a4paper,oneside]{report}
\usepackage{../style/consulting}
\hypersetup{pdftitle={Rapport d'Écarts — Time'Eats PMO}}
\pagestyle{fancy}\fancyhf{}
\renewcommand{\headrulewidth}{0pt}
\fancyhead[L]{\begin{tikzpicture}[remember picture,overlay]
  \draw[cnNavy,line width=0.5pt](0,0)--(\textwidth,0);
\end{tikzpicture}\vspace{2pt}\timeeatslogo\hspace{4pt}\small\textcolor{cnSlate}{Time'Eats \textbf{·} Rapport d'Écarts — \textit{[Nom projet]}}}
\fancyhead[R]{\small\textcolor{cnSlate}{\thepage}}
\fancyfoot[C]{\tiny\textcolor{cnSlate}{Document confidentiel — Cellule PMO — CESI MAALSI 2026}}

\newcommand{\ragV}{\cellcolor{cnGreen!20}\textbf{\textcolor{cnGreen}{Vert}}}
\newcommand{\ragA}{\cellcolor{cnOrange!20}\textbf{\textcolor{cnOrange}{Ambre}}}
\newcommand{\ragR}{\cellcolor{cnRed!20}\textbf{\textcolor{cnRed}{Rouge}}}

\begin{document}\small

\begin{center}
{\LARGE\bfseries\color{cnNavy} Rapport d'Écarts}\\[4pt]
\textcolor{cnOrange}{\rule{10cm}{1pt}}\\[4pt]
{\large\color{cnSlate} Phase 4 — Maîtrise \quad|\quad Projet : \textit{[Nom du projet]}}
\end{center}

\vspace{4pt}
\begin{alert}
Ce rapport est déclenché automatiquement dès qu'un indicateur passe en statut \textbf{Rouge} (IPC $<$ 0.85, IPD $<$ 0.85, retard $>$ 15j, ou risque critique actif). Il doit être transmis au DSI sous \textbf{48h}.
\end{alert}

\vspace{4pt}
\renewcommand{\arraystretch}{1.4}
\begin{tabularx}{\textwidth}{L{3.5cm} Y L{3.5cm} Y}
\toprule
\rowcolor{cnNavy}\th{Champ} & \th{Valeur} & \th{Champ} & \th{Valeur} \\
\midrule
Projet              & \textit{[Nom]} & N° rapport       & \textit{[ER-AAAA-XX]} \\
\rowcolor{cnIce}
Date de l'écart     & \textit{[JJ/MM/AAAA]} & CP          & \textit{[Prénom NOM]} \\
Type d'écart        & $\square$ Coût \quad $\square$ Délai \quad $\square$ Qualité \quad $\square$ Risque & Sévérité & \textit{[\ragR / \ragA]} \\
\bottomrule
\end{tabularx}

\section*{1 — Description de l'écart}

\textit{[Décrire précisément l'écart constaté : valeur attendue vs valeur réelle, magnitude de la dérive, depuis quand.]}

\section*{2 — Analyse des causes racines}

\renewcommand{\arraystretch}{1.3}
\begin{tabularx}{\textwidth}{C{0.6cm} L{5cm} Y C{3cm}}
\toprule
\rowcolor{cnNavy}\th{\#} & \th{Cause} & \th{Description} & \th{Type} \\
\midrule
1 & \textit{[Cause principale]} & \textit{[Explication]} & \textit{[Interne / Externe]} \\
\rowcolor{cnIce}
2 & \textit{[Cause secondaire]} & \textit{[Explication]} & \textit{[Interne / Externe]} \\
\bottomrule
\end{tabularx}

\section*{3 — Impacts prévisionnels}

\renewcommand{\arraystretch}{1.3}
\begin{tabularx}{\textwidth}{L{3cm} Y Y}
\toprule
\rowcolor{cnNavy}\th{Axe} & \th{Impact si non corrigé} & \th{Impact après correction} \\
\midrule
Délai   & \textit{[ex. +X semaines sur Jn]} & \textit{[ex. +X jours récupérables]} \\
\rowcolor{cnIce}
Budget  & \textit{[ex. +X K€ dépassement]} & \textit{[ex. dans réserve]} \\
Qualité & \textit{[ex. scope réduit]}       & \textit{[ex. fonctions maintenues]} \\
\bottomrule
\end{tabularx}

\section*{4 — Plan d'action correctif}

\renewcommand{\arraystretch}{1.3}
\begin{tabularx}{\textwidth}{C{0.6cm} Y L{3.5cm} C{2.5cm} C{2cm}}
\toprule
\rowcolor{cnNavy}\th{\#} & \th{Action corrective} & \th{Responsable} & \th{Échéance} & \th{Statut} \\
\midrule
A1 & \textit{[Action 1]} & \textit{[Prénom NOM]} & \textit{[JJ/MM]} & \textit{[À faire]} \\
\rowcolor{cnIce}
A2 & \textit{[Action 2]} & \textit{[Prénom NOM]} & \textit{[JJ/MM]} & \textit{[En cours]} \\
A3 & \textit{[Action 3]} & \textit{[Prénom NOM]} & \textit{[JJ/MM]} & \textit{[À faire]} \\
\bottomrule
\end{tabularx}

\section*{5 — Décisions demandées}

\begin{itemize}
  \item \textit{[Décision requise de la DSI : ex. autorisation d'utiliser la réserve]}
  \item \textit{[Décision requise de la DG si écart $>$ seuil]}
\end{itemize}

\vspace{8pt}
\renewcommand{\arraystretch}{2}
\begin{tabularx}{\textwidth}{Y Y}
\toprule
\rowcolor{cnNavy}\th{Validé par CP} & \th{Lu et approuvé DSI} \\
\midrule
\textit{Nom : \hfill Date :} & \textit{Nom : \hfill Date :} \\
\bottomrule
\end{tabularx}

\end{document}
